% ===== 原文 pp.24–25 =====
% tr03c：pp.24–25。p.24 首两条 bullet 承接 tr03b 末的 itemize（文件边界不能切断环境，
% 故此处另开一个新的 itemize，视觉上仍是同一列表）。
% [原文 p.24]
\begin{itemize}[leftmargin=1.6em,itemsep=2pt,topsep=2pt]
\item 对 $R-\cP$ 断开、从 $\beta_1$ 到 $\beta_1$ 的简单曲线求和时，被加函数是 $\cD$，且
\[
\frac{\partial\cD}{\partial L_1}=H(\ell_{\alpha_1}+\ell_{\alpha_2},L_1),
\]
它是多重曲线长度的函数；所得积分为
\[
\int_0^\infty\int_0^\infty\sum_{\substack{g_1+g_2=g\\ I_1\coprod I_2=\{2,\ldots,n\}}}
H(x+y,L_1)V_{g_1}(x,L_{I_1})V_{g_2}(y,L_{I_2})\,x y\,\dd x\,\dd y.
\]
\item 对 $j\neq1$、从 $\beta_1$ 到 $\beta_j$ 的简单曲线求和时，被加函数是 $\cR$，且
\[
\frac{\partial\cR}{\partial L_1}=\frac{1}{2}H(\ell_\alpha,L_1+L_j)+\frac{1}{2}H(\ell_\alpha,L_1-L_j),
\]
它是多重曲线加权长度的一些函数之和；所得积分为
\[
\int_0^\infty\frac{1}{2}\sum_{j=2}^n(H(x,L_1+L_j)+H(x,L_1-L_j))
V_g(x,L_2,\ldots,\widehat{L_j},\ldots,L_n)\,x\,\dd x,
\]
其中 $V_g$ 的自变量表中略去了 $L_j$。
\end{itemize}

与第 1 讲末尾的结果相比较——体积递推至此建立。\emph{体积函数 $V_{g,n}(L)$ 由递推确定。}

$V_{g,n}(L)$ 是怎样的函数？递推涉及两个初等积分（$H$ 的定义见前面的公式 (1)），
\[
\int_0^\infty x^{2j+1}H(x,t)\,\dd x\qquad\text{与}\qquad
\int_0^\infty\int_0^\infty x^{2j+1}y^{2k+1}H(x+y,t)\,\dd x\,\dd y.
\]
由直接计算（见 [Mir07a, 公式 (6.2) 与引理 6.2]），每个积分都是 $t^2$ 的多项式，其每个系数是阶乘之积与 Riemann zeta 函数在某个非负偶整数处取值之积，每个系数都是 $\pi$ 的适当幂的正有理倍数。第一个多项式在 $t^2$ 中的次数为 $j+1$，第二个的次数为 $i+j+2$。定理 4 的第一部分至此建立。第二部分是下一讲的主题。

% [Tutor Note: p.24 末"第二个的次数为 $i+j+2$"里的 $i$ 在全文中没有出现：紧挨着的公式
%  两个指数是 $x^{2j+1}y^{2k+1}$，指标只有 $j,k$，正确次数应为 $j+k+2$（原文笔误；
%  已用 700 dpi 裁图复核，该行确印作 $i+j+2$）。核对办法：$t$ 很大时当 $x+y<t$ 有
%  $H(x+y,t)\to1$，于是 $\iint x^{2j+1}y^{2k+1}H\approx\iint_{x+y<t}x^{2j+1}y^{2k+1}
%  =\frac{(2j+1)!(2k+1)!}{(2j+2k+4)!}\,t^{2j+2k+4}$，即 $t^2$ 的次数为 $j+k+2$。
%  数值复核（200 dpi/700 dpi 页图无关，纯积分计算）：$j=k=1$ 时该积分在 $t=20,40$ 的
%  $\log$–$\log$ 斜率约 $7.1$，随 $t$ 增大趋近 $2(1+1+2)=8$；而第一个积分在 $j=1$ 时
%  斜率约 $3.8$，趋近 $2(j+1)=4$。故 $j+k+2$ 与渐近一致，$i+j+2$ 不能成立。]

\para{第 4 讲：辛约化、$S^1$ 主丛与正规形式。}

本讲的目标是建立下面的定理。该公式与定理 3（WP kappa 方程 $\omega=2\pi^2\kappa_1$）相结合，即可说明体积多项式 $V_g(L)$ 的系数是 tautological 相交数。这一结果完成了定理 4 的证明。

% [原文 p.25]
\begin{thm}{10} 设 $d=\frac{1}{2}\dim_\R\Tg(L)=\dim_\C\Tgn$，则
\[
V_g(L)=\frac{1}{d!}\int_{\Tg(L)/\MCG}\omega^d_{\Tg(L)}
=\frac{1}{d!}\int_{\Tgn/\MCG}(\omega+\sum_{j=1}^n\frac{L_j^2}{2}\psi_j)^d.
\]
\end{thm}

证明本质上是在辛空间之间建立上同调等价——把辛约化、Duistermaat--Heckman 定理以及显式的几何结合起来，即得该公式。

本讲的讨论在相应的 Teichmüller 空间 $\cT$ 与开模空间 $\cT/\MCG$ 上展开。这些构造与增广（augmentation）构造相容。结果对相应的 Teichmüller 空间与紧化模空间成立。紧化模空间是 orbifold。以下结果都是系数取在 $\Q$ 中的上同调陈述；对流形而言，系数取在 $\Q$ 中的上同调论证一般对 orbifold 同样有效。此外，也可以绕过 orbifold 的麻烦：利用"紧化模空间具有流形的有限覆盖"这一一般结果，[BP00, Loo94]。

% [Tutor Note: 定理 10 的两个等号把"余体积"与"相交数"接起来：第一个积分是
%  $\Tg(L)/\MCG$ 上的 $\omega^d/d!$（$d$ 为复维数，故用 $d!$ 而非 $2^d d!$ 归一化）；
%  第二个积分把边界长度的影响写成 $\psi$ 类，$\omega+\sum_j\frac{L_j^2}{2}\psi_j$ 正是
%  Wolpert 的 $\dd\ell\wedge\dd\tau$ 公式在带边界情形下对 $\psi_j$ 的展开。可核对的维数锚点：
%  $\dim_\R\Tg(L)=6g-6+2n$，故 $d=3g-3+n=\dim_\C\Tgn$；取 $g=n=1$ 得 $d=1$，
%  定理 10 给出 $V(L)=\int_{\Mgn}(\omega+\frac{L^2}{2}\psi)$，而 p.32 用定理 3 的
%  $\omega=2\pi^2\kappa_1$ 把它写成 $V(L)=\int_{M_{1,1}}\left(2\pi^2\kappa_1+\frac{L^2}{2}\psi\right)$，
%  两式逐项一致，可用来检查 $\psi$ 类的系数是否漏了 $1/2$。]

\para{Teichmüller 空间。}我们考虑这个三元组。
\begin{itemize}[leftmargin=1.6em,itemsep=2pt,topsep=2pt]
\item $\Tg(L)$——亏格 $g$、带有指定长度 $L_1,\ldots,,L_n$ 的测地边界 $\beta_1,\ldots,\beta_n$ 的带标号双曲曲面所成之空间。双曲曲面可沿其测地边界作倍化（doubling），得到亏格更高的紧双曲曲面。相应地，$\Tg(L)$ 可看成 $\mathcal{T}_{2g+n-1}$ 中的一个轨迹（locus）。像 Teichmüller 空间的辛形式 $\omega$ 限制到该轨迹上，并在该轨迹上定义一个辛形式。带边曲面的裤子分解可以倍化，给出倍化曲面的裤子分解。倍化后的裤子分解以"包含测地线 $\beta_1,\ldots,\beta_n$ 且是对称的"为特征。对倍化分解应用 Fenchel--Nielsen 坐标与 $\dd\ell\wedge\dd\tau$ 公式，即得到轨迹 $\Tg(L)\subset\mathcal{T}_{2g+n-1}$ 的描述，并定义出辛形式 $\omega_{\Tg(L)}$。对带边曲面的任意裤子分解，该辛形式都表为 $\sum_j\dd\ell_j\wedge\dd\tau_j$（不含边界参数）。MCG 不变性是直接的。
\item $\what{\Tgn}$——亏格 $g$、有尖点（pointed）的测地边界 $\beta_1,\ldots,\beta_n$ 的带标号双曲曲面所成之空间——边界长度允许变化，并且每条边界上给定一个可变点。$\what{\Tgn}$ 的 $\R$ 维数比 $\Tg(L)$ 的 $\R$ 维数大 $2n$。辛约化要求 $\what{\Tgn}$ 上有一个辛形式，它在 $L$ 水平集上与 $\omega_{\Tg(L)}$ 等价，并且在边界上旋转这些点时不变。把 $\what{\Tgn}$ 描述成一个更高维的 Teichmüller 空间，即可给出这样的形式。为此，引入近乎紧的裤子（almost tight pants）——即两条有标号边界是穿刺点（长度为零）、第三条边界具有指定长度的裤子。每条曲面边界 $\beta_j$ 上都将粘上一个近乎紧的裤子。裤子的赤道反射（equatorial reflection）在裤子的边界上确定对称点；穿刺点的标号则唯一确定一个赤道边界点。
\end{itemize}

% [Tutor Note: p.25 第一项 "to define a symplectic from $\omega_{\Tg(L)}$" 的 "from" 应为
%  "form"（辛形式）；已用 700 dpi 裁图复核，确印作 "from"，译文按"辛形式"处理。
%  同页 "of prescribed lengths $L_1,\ldots,,L_n$" 处原文多印了一个逗号（700 dpi 复核确认），
%  属排印衍字，不影响含义，本译本照录。]
%
% 边界说明：p.25 末这一句（"An almost tight pants will be glued to each surface boundary β_j.
% The / pants equatorial reflection … uniquely determines an equatorial boundary point."）是
% p.25→p.26 的跨页句，按 TIMELINE 由 tr-03 负责写完，已写到此句的句号处。
% 紧随其后的 p.26 正文 "A standard model for a surface with pointed geodesic boundaries is
% given by gluing on almost tight pants …" 以及图 8、图 9 均由 tr-04 负责；该段仍属上面第 ②
% 条 bullet 的内容，tr-04 若要保持列表缩进，可用 \begin{itemize}\item[] … \end{itemize} 包裹。
